% Folder 74, batch 07 — archivists' pages 121–135.
% Pass: Opus 5.5 (claude-opus-5-5), 2026-10-03 — first pass, unchecked against the pages by a human.
%
% Skipped: p. 126 (blank coloured sheet, only a pencilled "(9p)" not in his hand).
% p. 122 is a printed USTL examination cover; only his sideways arithmetic is kept.
\documentclass[11pt,a4paper]{article}
\input{../preamble/grothendieck}

\folder{74}
\batch{7}
\pages{121}{135}
\foldertitle{Complexes cubiques : notes manuscrites (s.d.).}
\dating{[à partir de 1976]}
\watermark{Édition de démonstration}

\begin{document}

\section{Revêtements, torseurs et involutions $\sigma$}

\page{121}
\note{feuille de brouillon, notes éparses sans texte suivi ; l'argument vient d'avant le lot (les notations $V$, $\Gamma$, $\widetilde{E}_i$, $\sigma_x$ ne sont pas introduites ici). On donne les formules dans l'ordre approximatif de la page.}

\begin{tikzcd}[column sep=small, row sep=small, nodes={font=\scriptsize}]
 & 1 \arrow[d] & \\
V \arrow[r, "\varphi_i"] \arrow[dr, dashed, "\psi_i"'] & \Gamma \arrow[r, "p_i"] & E_i/2 = V \\
 & \widetilde{\Gamma} \arrow[u] \arrow[r, "\tilde{p}_i"] & \widetilde{E}_i \arrow[u] \\
 & \{\pm 1\} \arrow[u] & \pm 1 \arrow[u] \arrow[ul]
\end{tikzcd}

\note{lecture du diagramme incertaine : la flèche $p_i$ est surchargée, le but $E_i/2$ est en partie raturé, la flèche $\psi_i$ est en pointillé et porte la marque $\simeq$ ; sous $\{\pm1\}$ un « $0$ » \uncertain{} relié par un trait.}

$\widetilde{\Gamma} \to \widetilde{E}_i$ ; \qquad $g \in \widetilde{E}_1 = V$, \quad $\tilde{y} \in \widetilde{E}_2$.

\[
\begin{array}{lll}
x & y & z \\
x & y+u+w, & z+u+w \\
x_0 & y_0 + w & z_0 + w
\end{array}
\]
\struck{$x \ y \ z$} \quad \struck{$x \ \ z+u \ \ y+\ldots$}

\[
0 \to \widetilde{\Gamma} \to \widetilde{E}_1 \times \widetilde{E}_2 \to \qquad \widetilde{\Gamma} - \widetilde{E}_2
\]
\struck{$(g,y)$} $\mapsto \varepsilon$

$s \to t' \to u \to s'$ \qquad
$\widetilde{\Gamma} \simeq \widetilde{E}_1 \times \widetilde{E}_2 \to E_1 \times E_2$,
avec $\tilde{p}_3 : \widetilde{\Gamma} \to \widetilde{E}_3$, une flèche oblique $q_{12} : \widetilde{E}_1 \times \widetilde{E}_2 \to \widetilde{E}_3$, et $\widetilde{E}_3 \to E_3$ ;
\[
(x,y) \longmapsto \text{\struck{$x * y$}} \ (= -x\cdot y)
\]
\[
q_{12}(a\cdot\tilde{y}) = q_1\,\varepsilon(a)\,q_2\,\ldots \qquad 2^4/h = 2^2
\]
\note{la fin de la formule de $q_{12}$ est en abrégé : « $q_1 \varepsilon(a)\, q_2$ \uncertain{$\tilde{y}$} ».}

\struck{$\sigma_1 \ \sigma_2 \ \sigma_3$}

$\tau_1 \ \tau_2 \ \tau_3$
\drawing{un triangle de sommets $\widetilde{E}$ (en haut), $\widetilde{E}'$ (en bas à gauche), $\widetilde{E}_2$ \uncertain{} (à droite), les côtés marqués $\tau_3$, $\tau_1$, $\tau_2$.}
\[
\tau_2\tau_1 = \mathrm{id} \qquad \tau_3\tau_2 = \mathrm{id} \qquad \tau_1\tau_3 = \mathrm{id} \qquad \tau^2 = \mathrm{id}
\]

$\widetilde{G}$ sur $\widetilde{\Gamma}$ ; \qquad $\widetilde{\Gamma} \subset \widetilde{E}^3$
\[
\sigma_x\sigma_y\sigma_x^{-1}\sigma_y^{-1} = \sigma_{\sigma_x y}\,\sigma_y^{-1}
\]
\note{lecture de cette ligne incertaine ; sous l'indice $\sigma_x y$, un « $2$ » souligné.}
\[
\underbrace{\sigma_x\sigma_{x'}}_{\pi}\ \underbrace{\sigma_y\sigma_{y'}}_{\sigma_{\pi y}\sigma_{\pi y'}}\ \sigma_x\sigma_{x'}\ \sigma_y\sigma_{y'}
\]
\[
\pi_{x,x'}\,\pi_{y,y'} = \pi_{x,x'}\,\pi_{\pi_{x,x'}y,\ \ldots}
\]

\page{122}
\note{feuillet imprimé (copie d'examen) ; seul un calcul de sa main, écrit en travers, est transcrit.}
\[
12 \cdot 4^2 \cdot \text{\struck{$6$}}\,8 \cdot 9 \cdot 5 = 2^6 \cdot 3^3 \cdot 5
\]
\[
2^7 \cdot 3^4 \cdot 5 \qquad 128 \times 405 = 51840 \qquad /\,2\cdot 3
\]
\note{la multiplication posée $128 \times 405$ donne les produits partiels $640$ et $51200$ \uncertain{} et le total $51840$.}

\section{Graphes cubiques $\tau_\alpha$ et sous-structures}

\page{123}
\begin{itemize}
  \item \struck{$\tau$} $\tau_{-\infty}$ graphe vide
  \item $\tau_{-1}$ graphe \struck{\ill{}} triangle \qquad $c = 0 = 2^{-\infty}$
  \item $\tau_0$ \ \struck{\ill{}} \add{pavé} \qquad $c = 1 = 2^0$
  \item $\tau_1$ \qquad $c = 2 = 2^1$
  \item $\tau_2$ graphes cubiques \qquad $c = 4 = 2^2$
\end{itemize}

Cat.\ des graphes $\tau_\alpha$ munis d'une sous-structure $\tau_\beta$ ($\beta \leq \alpha$)
\[
\begin{array}{c|c|c|c|c}
\beta \backslash \alpha & -1 & 0 & 1 & 2 \\ \hline
-\infty & (\mathrm{Ens})_3 & (\mathrm{Ens})_{3,3} & (\mathrm{Ens})_6 & * \\
-1 & (\mathrm{Ens})_3 & (\mathrm{Ens})_3 \times (\mathrm{Ens})_2 & (\mathrm{Ens})_{2,2,2} & \text{Trialité} \\
0 & \emptyset & (\mathrm{Ens})_{3,3} & (\mathrm{Ens})_{3,3} & (\mathrm{Ens})_{3,3} \times (\mathrm{Ens})_{3,3} \\
1 & \emptyset & \emptyset & (\mathrm{Ens})_6 & (\mathrm{Ens})_6 \times (\mathrm{Ens})_2
\end{array}
\]
\note{plusieurs cases sont surchargées : sous $(\mathrm{Ens})_{2,2,2}$ un mot raturé \ill{} ; les indices $2,2,2$ et $3,3$ sont repassés à l'encre plus noire ; l'étiquette de la première ligne se lit $-\infty$ \uncertain{} (« $-\emptyset$ »), celle de la dernière est surchargée ($4$ ou $1$).}

Cas des systèmes $\tau_\alpha \subset \tau_\beta \subset \tau_\gamma$ ($\alpha < \beta < \gamma$ fixés)
\[
\begin{array}{c|c}
-1,\,0,\,1 & (\mathrm{Ens})_3 \\
-1,\,0,\,2 & (\mathrm{Ens})_3 \times (\mathrm{Ens})_3 \\
-1,\,1,\,2 & (\mathrm{Ens})_{2,2,2} \times (\mathrm{Ens})_2 \\
0,\,1,\,2 & (\mathrm{Ens})_{3,3} \times (\mathrm{Ens})_2 \\
-1,\,0,\,1,\,2 & (\mathrm{Ens})_3 \times (\mathrm{Ens})_2
\end{array}
\]
\note{ce second tableau est lu sur un aperçu sans agrandissement suffisant ; ses entrées sont toutes \uncertain{douteuses}.}

\page{124}
\[
\begin{array}{lll}
\sigma_{t_0} : & x \leftrightarrow x', \quad y \leftrightarrow y', \quad z \leftrightarrow z' \\
\sigma_t : & x \leftrightarrow y', \quad y \leftrightarrow z', \quad z \leftrightarrow x' \\
\sigma_{t'} : & x \leftrightarrow z', \quad y \leftrightarrow x', \quad z \leftrightarrow y'
\end{array}
\]
$\sigma_{t_0}\sigma_t$ : $x \to y \to z$, $x' \to z' \to y'$, $a \to s \to s'$ ; \quad
$\sigma_t\sigma_{t'}$ : $x \to y \to z$, \ldots, $a \to s \to s'$.
\[
\sigma_{t_0}\sigma_t = \sigma_t\sigma_{t'} = \sigma_{t'}\sigma_{t_0} = \pi_a, \qquad \pi^3 = \mathrm{id}
\]
\[
[\sigma_{t_0},\sigma_t] = [\sigma_t,\sigma_{t'}] = [\sigma_{t'},\sigma_{t_0}] = \pi^2 = \pi^{-1}
\]
\[
\sigma_t\sigma_{t_0} = \sigma_{t_0}\sigma_{t'} = \sigma_{t'}\sigma_t = \pi^{-1}
\]
\struck{$\sigma_{t_0}\sigma_t\sigma_{t'} = \sigma_t$ \ldots}

$\mathfrak{S}_I$ opère sur $S \smallsetminus X$ \ill{} fixé, \uncertain{3 avec pt fixe} \ldots\ \uncertain{trois orbites}, donc \ldots\ dans trois torseurs sous $\mathfrak{S}_I$ ; donc trois cas : \ill{}.
\note{le passage est rapide et en grande partie illisible ; seuls les termes mathématiques sont sûrs.}

\struck{$X$} \ $\downarrow$ \quad $Z$ cas : $9$ él. \quad $\downarrow$ \quad $E$ cas : $3$ él.

\drawing{trois ovales reliés, contenant $\xi_1\ \xi_2\ \xi_3$ / $\xi_4^*\ \xi_5^*\ \xi_6^*$ ; $\xi_4\ \xi_5\ \xi_6$ / $\xi_1'\ \xi_2'\ \xi_3'$ ; $\xi_{45}\ \xi_{56}\ \xi_{64}$ / $\xi_{12}\ \xi_{23}\ \xi_{31}$.}

II, III \quad deux cas : trois él.

$\mathfrak{S}_I \times \mathfrak{S}_J$ opère sur \ill{} \struck{\ill{}} fixe \ldots ; $\mathfrak{S}_I$ et $\mathfrak{S}_J$ opèrent \uncertain{sur} les pts fixes, donc \ldots\ $E = \coprod E_i$, de façon que les $E_i$ \uncertain{soient} des torseurs sous $\mathfrak{S}_I \times \mathfrak{S}_J$.

\page{125}
\note{page de calculs combinatoires sur les sommets $s,t,u$, $s',t',u'$, $x,y,z$, $x',y',z'$ d'un graphe ; on donne ce qui se lit.}

$(a\,b\,c)$ \qquad \struck{$x \in \widetilde{E}_a\ x$} \qquad $(s,t,u)$ \ $(s',t',u')$ \qquad $x \in \widetilde{E}_a \smallsetminus \{s,s'\}$

\[
\begin{array}{ll}
(s,t)\ (s',t') & (x,u)\ (x',u') \\
(t,u)\ (t',u') & (x,t)\ (x',t) \\
(u,s)\ (u',s') &
\end{array}
\]
$(u,x,y')$ \ $(u',x',y')$, \quad $y \in \widetilde{E}_b \smallsetminus \{t,t'\}$ ; \qquad
$(x,t',z')$ \ $(x',t,z)$, \quad $z \in \widetilde{E}_c \smallsetminus \{u,u'\}$
\[
\begin{array}{ll}
(y',u)\ (y,u') & (z',t')\ (z,t) \\
(y',s')\ (y,s) & (z',s)\ (z,s') \\
(y',x)\ (y,x') & (z',x)\ (z,x') \\
 & (z',y)\ (z,y')
\end{array}
\]

En marge, en haut à droite :
\[
\text{\struck{$3(1+2c)$}} \qquad \tfrac12(3+6c)(2+2c) = 3(1+2c)(1+c) = 3(1+3c+2c^2)
\]

\drawing{un graphe aux sommets $s'$, $s$, $u'$, $u$, $t$, $a$, $b$, $c$, $y'$.}

$s\ t\ u$ / $a\ b\ c$ ; \quad $x \mapsto y$ \ $(u\,x\,y)$ ; \quad $y \mapsto z$ \ $(s'\,y\,z')$ ; \quad $z \mapsto x'$ \ $(t\,z'\,x')$

\drawing{un hexagone orienté $z' \to x \to y' \to z \to x' \to y$, précédé d'une esquisse raturée ; à côté, des lignes $a - b - c \to a$ et leurs permutations, \uncertain{} pour la plupart.}

\[
\boxed{\sigma_t\sigma_s\sigma_u = \sigma_{t_0} \text{ sur } \widetilde{E}_a \smallsetminus \{s,s'\}}
\]
\[
\sigma_t\sigma_s = \sigma_{t_0}\sigma_u = \sigma_u\sigma_{t_0}
\]
\[
\sigma_t\sigma_s\sigma_t^{-1} = \sigma_{\sigma_t s} = \sigma_u
\]
\[
\sigma_t\sigma_s = \text{\struck{\ldots}}\ \sigma_u\sigma_t \quad \text{i.e.} \quad \underbrace{\sigma_u\sigma_t\sigma_s}_{=\ \sigma_{t_0} \text{ sur } S\smallsetminus\{t_0\}} = \sigma_t
\]
\marginal{short-cut !!}

$x\ x'$ / $s\ s'$ / $a,b,c$ / $t\ t'$ / $y\ y'$ / $u,u'$ / $z,z'$
\[
\sigma_t\sigma_y \neq \sigma_y\sigma_t : \ s \to u \to x \qquad\qquad \sigma_u\sigma_z = \sigma_z\sigma_u : \ s \to t \to x'
\]

\section{Graphes triangulaires $T_\nu$ et l'espace $V = \mathbb{F}_2^I$}

\page{127}
$*)$ \ $(S,A)$ graphe triangulaire $T_\nu$, avec \struck{$\nu \in \mathbb{N}$} $\nu = 0, 1 \ldots 2$, $c_\nu = 2^\nu \geq 1$.

$a \in S$, $I$ l'ens.\ des triangles de sommet $a$ (de cardinal $c_\nu + 1 = 2^\nu + 1$).

\struck{\ill{}} $G_I = \mathbb{F}_2^I$, il opère sur $S$ \struck{par les $\sigma_t$ ($t \in I$)} via $\sigma_t$ ($t \in I$).

\struck{$S_A = \{ s \in S \mid s \text{ lié à } a\}$}

$s \mapsto s'$ \ \struck{\ill{}} Pour $s \in S$ \uncertain{distinct de $a$} \ldots\ \ill{} deux triangles en $a$ \ldots\ $\{a,s\}$ \ldots
\note{deux lignes très rapides, en partie raturées.}

Donc $s \mapsto s'$ est une involution sans pt fixe de $X$, et \uncertain{$S \simeq X/\sigma$}. On peut regarder les $X_i$ ($i \in I$) comme des torseurs sous $\mathbb{F}_2$, \uncertain{donc} $\prod_{i\in I} X_i$ comme un torseur sous $V$, \uncertain{soit} $E$.

Soit $Y$ l'ens.\ \ldots\ $S \smallsetminus X \smallsetminus \{a\} = \{ x \in S \mid x \text{ non lié à } a\}$.

Pour $x \in Y$ et $i \in I$, $x$ est lié à un élt \uncertain{unique} de $X_i$, \uncertain{soit} $\varphi_i(x)$, donc
\[
\varphi(x) = (\varphi_i(x))_{i\in I} \in E = \prod X_i ,
\]
\struck{d'où une application \ldots} $\varphi : Y \to E$.

Comme $G$ opère \ldots\ \ill{} ; \struck{\ill{}} Soit $z_0$ l'élément \uncertain{non nul} de $\mathbb{F}_2$, $z$ \uncertain{l'élément diagonal} de $\mathbb{F}_2^I$.

\struck{sur l'application diagonale \ldots\ $\sigma_i$ ($i \in I$)}

$\xi_i : \mathbb{F}_2 \to V \simeq \mathbb{F}_2^I$ l'inj.\ canonique, $\rho$ l'op.\ naturelle de $V$ sur $E$ par
\[
\rho(\xi_i(z_0))\,(x_j)_j = (y_j)_j , \quad y_j = x_j \text{ si } j \neq i, \quad y_i = z_0 x_i \text{ sinon}.
\]
On aura \struck{$\sigma_i(x) =$}

\page{128}
pour $x \in E = \prod_{i\in I} X_i$, et \struck{$g = (g_i)_{i\in I} \in V$} $i_0 \in I$,
\[
(\sigma_{i_0})_E = \prod_{i \in I \smallsetminus i_0} \rho(\xi_i(z_0)) = \rho\Bigl(\sum_{i\in I\smallsetminus i_0} \xi_i(z_0)\Bigr) = \rho(z)\,\rho(\xi_{i_0}(z_0))
\]
\note{sous la somme, une accolade : « $= z - \xi_{i_0}(z_0)$ » \uncertain{}.}

Donc si \ldots\ $V \xrightarrow{\ \eta\ } \mathbb{F}_2$ \uncertain{l'augmentation} ($\delta : \mathbb{F}_2 \to V$ diag.), on aura, pour $g \in V$ :
\[
(\sigma_g)_E = \rho\bigl(g + \delta\eta(g)\bigr)
\]

Considérons l'application
\[
V \xrightarrow{\ u\ } V, \qquad u(x) = x + \delta\eta(x)
\]
Son noyau est formé des $x$ tels que
\[
x = \delta\eta x
\]
\struck{\ill{}} i.e.\ c'est un \uncertain{sous-espace} \ldots\ $z = \delta(z_0)$, et il vient à voir \ldots
\[
\text{\struck{$z$}} = \delta\eta z \quad \text{i.e.} \quad \text{\struck{$z_0 = \eta\delta z_0$}} \quad (\text{comme } z = \delta z_0)
\]
quand \ldots
\[
z_0 = \eta z \quad \text{i.e.} \quad \eta z \neq 0
\]
: c'est le cas ssi $I$ est de cardinal impair. Or a cardinal $1 + c_\nu = 1 + 2^\nu$, qui est impair sauf pour $\nu = 0$ (cas du point). \struck{Donc}

\uncertain{dans ce cas}, si $\nu \neq 0$, $u$ est injectif donc iso, et \struck{\ill{}} \add{les op.}\ de $V$ sur $E$ via $\rho$ ou via $\sigma$ se déterminent mutuellement. Or l'application $\varphi$ commute

\page{129}
aux op.\ $\sigma$, on trouve que l'appl.\ $\varphi$ de $Y$ dans $E$ est \struck{\ldots} (\uncertain{donc injective}, \uncertain{card} $Y = \ldots$) sous $V$ et une \ldots\ \uncertain{donc} \ldots

$V$ \uncertain{lt} action. Or $(4c_\nu = 4$ \struck{\ill{}} $\mathrm{card}\, E = \ldots$
\[
\text{\struck{\ill{}}} = 2^{c_\nu + 1} = 2^2 = 4,
\]
donc \struck{\ldots} $\varphi$ est bijectif.

b) Si $\nu \neq 0$ i.e.\ $\nu = 1$ ou $2$, alors l'image \struck{\ill{}} $u$ \ldots
\[
\text{\struck{Im}}\ \delta(\mathbb{F}_2) = \{0, z\}, \quad \text{comme} \ \ldots
\]
\struck{(\ldots)}

\uncertain{noyau} \ldots\ $\operatorname{Ker} \eta = V$ \ \struck{\ill{} \ldots, $Y \hookrightarrow V$ \ldots}

\struck{\ill{}}
\note{l'argument s'interrompt ici ; le reste de la page est blanc.}

\section{Graphes pseudotrialitaires en termes de triangles}

\page{130}
\note{le titre de la section est de sa main, en tête de page.}

$X$ ens.\ (« sommets ») ; \quad $T \subset \mathfrak{P}_3(X)$ (triangles)

\emph{Axiomes} \quad
a) $\underline{t}, \underline{t}' \in T \Rightarrow \mathrm{card}(\underline{t} \cap \underline{t}') \neq 2$,
i.e.\ si $\underline{t}, \underline{t}'$ ont deux pts en commun, ils sont égaux ; ou encore : toute arête contenue dans un triangle \emph{unique}.

Déf : $\underline{a} \in \mathfrak{P}_2(X)$ est une arête si $\exists\, \underline{t} \in T$, $\underline{a} \subset \underline{t}$.

b) Soit $\underline{t} \in \mathfrak{P}_3(X)$, pour que $\underline{t} \in T$, il (faut et) suffit que $\forall\, \underline{a} \in \mathfrak{P}_2(X)$, $\underline{a} \subset \underline{t}$, on ait $\underline{a}$ est une arête. On aura \ldots\ [lorsque $a,b,c \in X$ distincts, $\underline{t}, \underline{t}', \underline{t}'' \in T$ tels que $a,b \in \underline{t}$, $b,c \in \underline{t}'$, $c,a \in \underline{t}''$, alors $\{a,b,c\} \in T$ (\uncertain{et} $\underline{t} = \underline{t}' = \underline{t}'' = \{a,b,c\}$)].

c) $\forall\, \underline{t} \in T$ et $a \in X$, $\exists\, \underline{t}' \in T$ tel que $a \in \underline{t}'$, $\underline{t} \cap \underline{t}' \neq \emptyset$ [et si $a \notin \underline{t}$, $\underline{t}'$ est unique \ldots]

Pour que $X$ soit une \uncertain{conf.}, il faut et il suffit qu'on ait

d) $\forall\, a \in X$, l'ens.\ des $\underline{t} \in T$ tels que $a \in \underline{t}$ est \struck{un triangle} \uncertain{de cardinal $3$}.

Soit $X$ une \uncertain{conf.}, \struck{$X$} $X'$ l'ens.\ des \struck{triangles} \uncertain{de ses droites}, $R$ la relation d'incidence. Donc $\forall\, x' \in X'$, $R\{x'\} \in \mathfrak{P}_3(X)$, et $x' \mapsto R\{x'\}$ de $X'$ dans $\mathfrak{P}_3(X)$ est injectif\ldots\ Mais aussi, $\forall\, x \in X$, $R\{x\} \in \mathfrak{P}_3(X')$, et $x \mapsto R\{x\}$ de $X$ dans $\mathfrak{P}_3(X')$ est injectif.

Soit alors $X$ une \uncertain{\ill{}}. Montrons que $X \subset \mathfrak{P}_3(X')$ satisfait aux conditions a) b) c) d) précédentes.
\note{ici et à la page suivante, $\mathfrak{P}_k(X)$ rend son $\mathcal{P}$ gothique (ensemble des parties à $k$ éléments) ; les triangles et arêtes sont soulignés sur la page.}

\page{131}
Vérifions d) : $\forall\, a' \in X'$, l'ens.\ des $x \in X$ incidents à $a'$ est bien trois\ldots

a) Soient $a, b \in X$, \add{$a \neq b$}, considérons les $\underline{t} \in X'$ incidents à $a$ et $b$ ; il faut, je dis qu'il y en a $0$ ou $1$ (jamais $2$) : clair.

b) Soient $\underline{t}, \underline{t}', \underline{t}'' \in T = X'$ \add{il.\ distincts} et $a,b,c \in X$ tels que $\underline{t}, \underline{t}'$ incidents à $a$, $\underline{t}', \underline{t}''$ inc.\ à $b$, $\underline{t}'', \underline{t}$ inc.\ à $c$, prouvons que $a = b = c$, donc que $\{\underline{t}, \underline{t}', \underline{t}''\}$ \uncertain{\ill{}} : un \ldots\ point\ldots\ En effet, si on avait $b \neq a$, le triangle $\underline{t}''$ ne pourrait venir rencontrer $\underline{t}$, car $b$ ne serait pas lié à aucun autre sommet de $\underline{t}$ que $a$.

\drawing{triangle de sommet $a$, côtés $\underline{t}$ et $\underline{t}'$, avec $b$ \uncertain{?} et $c$.}

c) $\forall\, a \in X$ et $\underline{t} \in T = X'$ \add{($a$ non incident à $\underline{t}$)}, $\exists\, b \in X$ incident à $\underline{t}$, et tel que $a, b$ soient incidents à un même $\underline{t}' \in T = X'$.

\drawing{un point $a$ et un triangle.}

\section{Diagrammes $A_n$ dans les graphes cubiques}

\page{132}
\note{page de dessins.}

\drawing{trois motifs de points ; deux triangles se touchant par un sommet ; un petit cube dont les sommets portent $\lambda_{25}, \lambda_{16}, \lambda_{13}, \lambda_{24}, \lambda_{14}, \lambda_{23}, \lambda_{26}, \lambda_{15}$, avec un trièdre d'axes ; un grand cube aux sommets alternativement pleins et creux ; plus bas, un polyèdre en treillis, un prisme et un carré.}

$\lambda_{12}$ \quad
$\left\{\begin{array}{l}\lambda_{46},\ \lambda_{35} \\ \lambda_{34},\ \lambda_{56} \\ \lambda_{45},\ \lambda_{36}\end{array}\right.$
\qquad $e_1 + e_2 \,/\, e_2 + e_3 \,/\, e_3 + e_1$

cube $\longleftrightarrow$ \uncertain{cube ponctué}
\[
(\mathrm{Ens})_{2,2,2} \qquad (\mathrm{Ens})_2 \times (\mathrm{Ens})_4
\]
\[
\mathfrak{S}_3 \cdot (\pm 1)^3 \simeq \mathfrak{S}_2 \times \mathfrak{S}_4
\]

\page{133}
\struck{$A_1(S_i)$}
\marginal{\uncertain{dilemme pt, triangle -- pas biseau ni tripode}}

diagramme $A_2 \simeq$ triangle ponctué \struck{\ill{}}
\drawing{un triangle de sommet $a$, deux côtés en pointillé, la base en trait épais.}

Cat.\ des graphes cubiques avec diagramme $A_2 \cong$ catégorie des ensembles de cardinal $8$ munis d'une involution ss pt fixe $\simeq$ catégorie des $4$-\uncertain{cubes} $\simeq$ catégorie des torseurs sous
\[
\mathfrak{S}_4 \cdot \text{\struck{$\mathbb{Z}$}}(\pm 1)^4 \simeq W_{B_4} \quad (\text{groupe de Weyl de } SO(9))
\]

diagramme $A_3$
\[
\zeta_1 \overset{\lambda_{12}}{\text{------}} \zeta_1^*
\]
\marginal{\uncertain{dilemme pt, biseau, pas triangle ni tripode}}
$\simeq$ doublet $+$ élément de la demi-étoile associé

$\simeq$ biseau $B$ $+$ él.\ marqué de $I = B/\sigma_B$ $+$ $2^{\text{ème}}$ él.\ marqué

($\zeta_1, \lambda_{12}, \zeta_1^*$) él.\ typique

Cat.\ des graphes cubiques avec diagramme $A_3 \cong$ \struck{cat.\ des couples} $(\mathrm{Ens})_2 \times (\mathrm{Ens})_4$
\drawing{un cube en perspective, une face en traits épais, des arêtes en pointillé.}

diagramme $A_4$

$\simeq$ pavé $+$ triade et triangle (sécants)

\struck{\ill{}} $\simeq$ hexagone $+$ \add{clou} \uncertain{centré}

$\cong (\mathrm{Ens})_2 \times (\mathrm{Ens})_3$
\drawing{un hexagone dont la moitié inférieure est en trait épais, avec un « W » inscrit ; trois rangées de points.}

diagramme $A_5 \simeq$ (\struck{pavé} \add{carré} ponctué)

$\simeq$ \struck{\ill{}} \add{hex.\ ponctué}

$\simeq$ triade épinglée

$\cong \mathrm{Ens}_{2,\text{\struck{4}}} \times \mathrm{Ens}_3$
\note{l'indice de $\mathrm{Ens}$ est surchargé : un $4$ repris en $2$ \uncertain{}.}
\drawing{un hexagone avec un point central $a$, les arêtes marquées $\xi_1$, $\xi_1^*$ \uncertain{} ; un dessin entièrement raturé ; un hexagone contenant un losange centré ; un hexagone divisé par des diagonales ; quelques points et flèches.}

\page{134}
Prismes triangulaires \struck{\ill{}} $\simeq$ pavé $+$ triangle
\[
\cong (\mathrm{Ens})_2 \times (\mathrm{Ens})_3 \times (\mathrm{Ens})_3
\]

\drawing{un cube raturé ; un prisme surmonté d'une pyramide, avec les points $a', a, a''$ et $b', b, b''$ ; à droite une figure analogue développée ; une longue flèche vers le mot « bicarré ».}

les triangles sur $(a'b')$ et $(a''b'')$ \ldots\ $=$ sommets ; si $t$ est le \add{$3^{\text{e}}$ sommet du} triangle de base $(a,b)$, \uncertain{le cube de} troisième sommet \add{$u$} du triangle de base $(ab)$ est le carré $(a\,b\,b'\,a')$ et du carré $(a\,b\,b''\,a'')$. Dans les deux cas \ldots\ distincts, on ne \ldots\ pas \ldots\ $=$ pavé : ils engendrent \ldots\ On peut dire aussi \ldots\ un $A_5$ (contenu dans un \ldots) minimum, en fait deux, un hexagone « ponctué » de façon unique \ldots\ du prisme : la triade $(a, b', b'')$ \ldots\ et il faut choisir un pt $b$ lié aux trois pts de cette triade, i.e.\ un W \ldots : la triade \uncertain{apposée}. Cas des deux en tout
\[
\tfrac12(720 \times 6 \times 3) = 2^4 \cdot 3^4 \cdot 5 = N : 2^3 \ \text{---}
\]
en effet, la cat.\ des \struck{\ill{}} graphes cubiques avec bicarré est
\[
\simeq \underbrace{(\mathrm{Ens})_2 \times (\mathrm{Ens})_2}_{\text{bicarrés}} \times \underbrace{(\mathrm{Ens})_2}_{\text{\uncertain{complétant l'enveloppe quasi-trialitaire}}} \simeq \mathrm{Tors}\bigl((\pm1)^3\bigr)
\]
\note{la moitié gauche de la page, entre les dessins, est très rapide ; la ponctuation des passages illisibles est rendue par « \ldots ».}

\page{135}
\drawing{une bande de trois carrés de sommets $a, a', a'', a'''$ en haut et $b, b', b'', b'''$ en bas.}
n'existe pas dans un graphe cubique, car on trouve que les \struck{triangles de base} \add{$3^{\text{e}}$ sommet} des triangles de base $(a''', b''')$ égale celui du triangle de base $(a'', b'')$ et $(a', b')$, ce qui donne \uncertain{égaux}, absurde \ldots\ d'après \ldots\ des \add{carrés} \ldots

Conséquence : pas de prismes à base pentagonale et hexagonale. Par contre \ill{} prismes à base carrée, i.e.\ cubes.

\marginal{NB. Carré définit pt, paire de tripodes \uncertain{} ; non biseau ni triangle\ldots}

Carrés de centre un pt donné $a$ : d'abord le \uncertain{domaine} des graphes cubiques \uncertain{à} pt marqué $a$, \uncertain{équivaut} : celle d'un ens.\ $I$ de card.\ $5$ et d'un torseur $T$ sous $V_I = (\mathbb{F}_2^{\,I})'$ (de dim.\ $4$ sur $\mathbb{F}_2$) \add{engendré} par des $e_i$ ($i \in I$) avec la seule relation $\sum e_i = 0$. L'ens.\ des pts non liés à $a$ est $T$. Les donnée\supplied{s} \ldots\ carrés dans les \uncertain{demi-espaces} \ldots\ (\uncertain{espaces} de dim.\ $2$, donc l'espace des translations engendré par $2$ parmi les $e_i$, \struck{\ill{}} \add{corr.} : une partie $I'$ de $I$ de cardinal $2$. Si $V_{(I')}$ est le s.-espace correspondant, les carrés \uncertain{en} question correspondent aux pts de $T/V_{(I')}$, \uncertain{lesquels} sont \uncertain{en} \uncertain{bijection} avec les \uncertain{quotients} du groupe \uncertain{fondamental} de $V_I$ : $V_{(I')} \simeq \mathbb{F}_2^{\,I'}$ [Donc la cat.\ des graphes cubiques $+$ carrés équivaut, celle des \struck{couples} \add{triples} \struck{(\ldots)} ($I'$ ens.\ de card.\ $2$, $T'$ torseur sous $\mathbb{F}_2^{\,I'}$, $I''$ ($= I \smallsetminus I'$) ens.\ de card.\ $3$) i.e.
\[
\text{\struck{$(\mathrm{Ens})$}}_{2,2} \times \mathrm{Ens}_3 \simeq \text{Carré} \times (\mathrm{Ens})_3
\]
\note{l'argument continue au-delà du lot : la phrase ouverte par le crochet n'est pas refermée sur la page.}

\end{document}
